\documentclass[12pt]{article}
\usepackage[utf8]{inputenc}
\usepackage[T1]{fontenc}
\usepackage{hyperref}
\title{Vai Language Resources for K-12 Curriculum Development}
\author{Compiled by Claude Sonnet 5 (claude-sonnet-5) for Andr\'as Kornai}
\date{}
\begin{document}
\maketitle
\section{Introduction}
Vai (ISO 639-3 vai, Glottocode vaii1241) is a Mande language of Liberia and Sierra Leone. Wikipedia lists roughly 104,000 speakers in Liberia and 15,500 in Sierra Leone. It has its own syllabary, traditionally credited to Momolu Duwalu Bukele around 1833; one source attributes it differently.
\section{Primary Educational Resources}
My search found no Vai primer or textbook online. It mentioned that the Institute of Liberian Languages has published Vai materials; I did not see a catalogue.
\section{Linguistic and Reference Materials}
\begin{itemize}
\item Vai language, Wikipedia: \url{https://en.wikipedia.org/wiki/Vai_language}
\item Vai syllabary, Wikipedia: \url{https://en.wikipedia.org/wiki/Vai_syllabary}
\item Vai Unicode block (U+A500--U+A63F): \url{https://en.wikipedia.org/wiki/Vai_(Unicode_block)}
\item Unicode proposal chart for Vai: \url{https://www.unicode.org/L2/L2005/05053r-vai-chart.pdf}
\end{itemize}
\section{Religious and Community Resources}
The Unicode documentation reports that a Qur'an translation and a New Testament translation exist in Vai. I did not find a URL for either.
\section{Digital and Audio Resources}
The script is encoded in Unicode (version 5.1, 2008), which makes digital text possible. No audio resources were found.
\section{Curriculum Development Context}
The syllabary and Unicode support are real assets. Teaching materials, dictionaries and school programmes were not located and need further search, for example through the Institute of Liberian Languages.
\end{document}